\section{Historical motivation and provenance of the distance endpoint}
\label{app:time-current}

The material in this appendix records the theoretical route by which
the distance-compatible endpoint was discovered. None of the bridge or
triangulation claims requires this interpretation.

The investigation began from a redshift-dependent time-current
framework in which a global response $\Gamma(z)$ rescales the
background expansion rate, recovering the reference expansion in the
standard limit $\Gamma(z)\to 1$. The M24--M28 development mutated
that response against DESI DR2, Pantheon+SH0ES, and compressed-CMB
likelihoods, passing through early-window, ruler-mechanism, and
Xi-mapping-repair stages documented in
Appendix~\ref{app:mutation-history}. The distance endpoint~R used in
the present paper descends from that search lineage as a smooth
low-redshift control history; its fitted time-current parameters are
a record of discovery, not inputs to the reconstruction.

The perturbative-completion program attempted to extend the fitted
background with derived perturbation dynamics (direct variation,
symmetry arguments, effective-field routes, lapse
reinterpretations). Each route failed to produce an admissible
extension, and the homogeneous-ansatz route was refused as an
unconstrained knob. Those failures are evidence fixing the scope of
the present paper to background geometry: they are why chronometer,
growth, lensing, time-delay, and siren sectors are recorded as not
predicted without a covariant completion.

No result in Sections~\ref{sec:bridge} or \ref{sec:triangulation}
assumes the lapse/current variables, the $\Gamma(z)$ law, or any
property of the time-current dynamics beyond the frozen numerical
tables of~R. A reader who rejects the entire time-current idea can
still evaluate the empirical comparison, because the comparison
operates only on the frozen histories.
